% !TEX root = ../../main.tex

\section{系统设计目标与原则}

在完成系统需求分析后，需要进一步对 HearBridge 系统的总体架构、功能模块、数据结构和关键业务协作方式进行设计。
系统设计的目标是将第三章中提出的普通用户训练识别需求和管理员资源模型维护需求转化为可实现的系统结构，使移动端、后台管理端、服务端和手势识别服务能够形成清晰的协作关系。

HearBridge 系统的设计主要遵循以下原则。

第一，前后端分离原则。
系统将移动端、后台管理端和服务端进行分离，移动端与后台管理端主要负责界面展示和用户交互，Spring Boot 服务端负责业务逻辑、数据管理和跨服务编排。
前后端之间通过统一接口进行通信，有利于降低界面层和业务层之间的耦合度。

第二，多端协同原则。
系统同时包含 HarmonyOS 移动端、Vue Web 管理端、Spring Boot 服务端和 Python 手势识别服务，各端承担不同职责。
移动端面向普通用户提供训练和识别入口，管理端面向管理员提供资源和模型维护入口，服务端负责统一业务管理，识别服务负责手势识别相关能力。
通过多端协同，可以使系统功能边界更加清晰。

第三，识别服务独立原则。
手势识别、句子视频识别、样本处理和模型加载等能力由 Python 手势识别服务单独承担，而不直接放在移动端或主业务服务中。
这样既便于利用 Python 生态中的视觉识别和深度学习工具，也便于后续独立调整识别策略和模型版本。

第四，资源与模型可维护原则。
系统不仅需要面向普通用户提供训练和识别功能，还需要面向管理员提供手势分类、手势资源、样本数据和模型版本等维护能力。
通过后台管理端和服务端的数据管理功能，可以为后续扩充训练资源、优化识别模型和发布新模型版本提供基础。

第五，有限场景可控原则。
系统当前定位为听障辅助训练与受限场景句子视频识别原型，不追求开放域连续手语翻译。
系统设计中保留识别结果、语义修正结果和模型版本等信息，使其能够在有限词表和受控场景下形成较清晰的功能闭环。
